\section{Ingest and Indexing}
\label{sec:ingest}

This section covers how process events enter \vannak{}, how they are validated
and deduplicated, how they reduce into shard-local state, the query primitives
over that state, and how data-individual metadata events are captured.

\subsection{Process Event Ingest}

Process events enter as \texttt{PipelineEvent} values and are submitted to a
shard-local \texttt{HotIndex} through \texttt{HotIndex::ingest}. The index is
intentionally single-threaded; the expected \sitas{} integration holds one
\texttt{HotIndex} per owning shard.

\subsection{Validation}

Before reduction, \texttt{HotIndex::ingest} calls
\texttt{PipelineEvent::validate}, which requires that the core identity fields
are non-empty: \texttt{event\_id}, \texttt{source\_id}, \texttt{tenant\_id},
\texttt{environment\_id}, \texttt{pipeline\_id}, \texttt{process\_definition\_id},
\texttt{process\_instance\_id}, and \texttt{timestamp}. A violation returns
\texttt{IngestError::EmptyField} naming the offending field.

\subsection{Deduplication}

The index deduplicates by \texttt{event\_id}. If an event with the same id has
already been ingested, the index increments its \texttt{duplicate\_events}
counter and returns \texttt{IngestOutcome::Duplicate} without modifying state.
This makes re-delivery of the same event safe and idempotent.

\begin{lstlisting}
pub fn ingest(&mut self, event: PipelineEvent) -> Result<IngestOutcome, IngestError> {
    event.validate()?;
    if self.events_by_id.contains_key(event.event_id()) {
        self.duplicate_events += 1;
        return Ok(IngestOutcome::Duplicate);
    }
    // ... reduce state and index ...
    Ok(IngestOutcome::Accepted)
}
\end{lstlisting}

\subsection{IngestOutcome}

\texttt{IngestOutcome} is a small enum with two variants: \texttt{Accepted} when
the event was reduced and indexed, and \texttt{Duplicate} when an event with the
same id was already present.

\subsection{Durga Mirroring}

The \texttt{durga} module mirrors \durga{}'s canonical monitoring event as
\texttt{DurgaProcessEvent}, with \texttt{DurgaStatus} and \texttt{DurgaEventType}
that convert into \vannak{}'s \texttt{EventStatus} and \texttt{EventKind}. The
method \texttt{into\_pipeline\_event} builds a \texttt{PipelineEvent},
synthesizing the event id from the source id, process instance id, and source
sequence, and copying optional fields such as activity, token, correlation,
process version, business key, error, and payload.

\subsection{HotIndex Structure}

\texttt{HotIndex} keeps the current process state plus several secondary
indexes, all keyed for ordered iteration:

\begin{description}
\item[\texttt{process\_instances}] current \texttt{ProcessInstanceState} per
  instance.
\item[\texttt{events\_by\_id}] full event store keyed by \texttt{event\_id}.
\item[\texttt{events\_by\_process}] event ids per process instance.
\item[\texttt{events\_by\_pipeline}] event ids per pipeline.
\item[\texttt{events\_by\_metadata\_ref}] event ids per metadata reference, for
  impact analysis.
\item[\texttt{duplicate\_events}, \texttt{rejected\_events}] counters surfaced
  in snapshots.
\end{description}

\subsection{Process State Reduction}

Each accepted event is folded into the instance's
\texttt{ProcessInstanceState::apply}. The reducer updates the coarse
\texttt{ProcessStatus}, tracks the current activity, records per-activity
\texttt{ActivityState}, and computes activity durations. When an activity
reaches a terminal state, the reducer subtracts its recorded entry time from the
event timestamp using \texttt{EventTimestamp::saturating\_duration\_until},
producing a duration in milliseconds. A \texttt{ProcessFailed} event increments
\texttt{retry\_count} and records the last error.

\begin{table}[htbp]
\centering
\begin{tabular}{ll}
\hline
Event kind & Resulting process status \\
\hline
\texttt{ProcessStarted} & \texttt{Active} \\
\texttt{ActivityEntered} & \texttt{Active} \\
\texttt{ActivityCompleted} & \texttt{Active} \\
\texttt{GatewayTaken} & \texttt{Active} \\
\texttt{ActivityCancelled} & \texttt{Cancelled} \\
\texttt{ProcessCompleted} & \texttt{Completed} \\
\texttt{ProcessFailed} & \texttt{Failed} \\
\hline
\end{tabular}
\caption{Process status transitions by event kind.}
\end{table}

\subsection{Query API Types}

The \texttt{query} module defines the typed requests and a single owned result
type. \texttt{QueryLimit} clamps to a minimum of 1 and defaults to 100.

\begin{description}
\item[\texttt{ProcessInstanceQuery}] fetch one instance snapshot by id.
\item[\texttt{PipelineQuery}] list instance snapshots for a pipeline.
\item[\texttt{EventQuery}] list events for a process instance.
\item[\texttt{ImpactQuery}] list events touching a metadata reference.
\item[\texttt{QueryResult}] either \texttt{ProcessInstances} or \texttt{Events}.
\end{description}

The matching \texttt{HotIndex} methods are \texttt{process\_instance},
\texttt{pipeline\_instances}, \texttt{events}, \texttt{impact}, and
\texttt{affected\_pipelines}. All return owned values; none expose borrowed
references into shard-local state.

\subsection{Data-Individual Metadata Events}

Provenance is captured with \texttt{DataIndividualMetadataEvent}. The event
carries identity (\texttt{DataIndividualId}, \texttt{MetadataEventId}), domain
placement (\texttt{DataIndividualShardId}), process context, a timestamp, a
\texttt{MetadataOperation}, and the passive and active metadata maps. A builder
attaches optional activity id, passive metadata, active metadata, and source
payload reference.

\subsubsection{Passive and Active Metadata}

\texttt{PassiveMetadata} captures what is observed at receive or create
boundaries (source system, format, size, checksum, received time).
\texttt{ActiveMetadata} captures what processing did (transformation plugin,
masking, validation result, enrichment source, field changes). Both are
insert-builders over \texttt{MetadataFieldName} to \texttt{MetadataValue}.

\subsubsection{MetadataOperation}

The operation classifies the provenance event: \texttt{Created},
\texttt{Received}, \texttt{Transformed}, \texttt{Masked}, \texttt{Validated},
\texttt{Enriched}, \texttt{Routed}, and \texttt{Persisted}. Several variants
carry structured detail, such as the plugin name and version on
\texttt{Transformed} or the pass flag on \texttt{Validated}.

\subsection{Idempotency Key}

Each metadata event derives a stable \texttt{IdempotencyKey} of the form
\texttt{data\_individual\_id:metadata\_event\_id}. This key drives outbox
deduplication and, downstream, the correlation id used for idempotent \ipto{}
writes.

\subsection{HotIndexSnapshot}

\texttt{HotIndex::snapshot} returns an owned \texttt{HotIndexSnapshot} with
counts only:

\begin{lstlisting}
pub struct HotIndexSnapshot {
    pub process_instance_count: usize,
    pub event_count: usize,
    pub pipeline_count: usize,
    pub metadata_ref_count: usize,
    pub duplicate_events: u64,
    pub rejected_events: u64,
}
\end{lstlisting}

\begin{vnote}
Snapshots are point-in-time copies of counters, not live views. Reading a
snapshot never blocks or borrows the index, which keeps observability cheap and
safe to expose across shard boundaries.
\end{vnote}

\endinput
